% ============================================================
% VII. CONCLUSIONS
% ============================================================

\section{Conclusions}

The analysis of \texttt{Muestra01.csv} showed that its dominant spectral content changes across the four $0.05$ s segments evaluated ($1200.00$, $1600.00$, $940.00$, and $1300.00$ Hz), indicating that it is a non-stationary signal within the recorded window.

Sampling the $100$ Hz sine experimentally confirmed the Nyquist theorem: at $70$ Hz ($f_s < 2f_{\max}$) the four samples obtained in $0.05$ s do not unambiguously represent the original signal and are consistent with aliasing, while at $500$ Hz and $1000$ Hz ($f_s \geq 2f_{\max}$) the quantized waveform faithfully follows the reference sine. Preventing aliasing in practice requires satisfying $f_s \geq 2f_{\max}$ and/or analog filtering of content above $f_s/2$ before sampling.

The FFT comparison of violin, drum, and cat showed three clearly separated dominant frequencies ($385.50$, $7596.10$, and $1715.08$ Hz respectively), confirming that frequency analysis does allow these three sound sources to be differentiated from one another. However, the drum's case — whose dominant peak reflects its attack transient rather than a low fundamental tone — shows that this differentiation does not always correspond to an intuitive musical interpretation and must be read carefully depending on whether the source is tonal or percussive in nature.

On ESP32-WROOM sampling precision, the captured data (from real hardware) shows that the single-threaded busy-loop mechanism is the only one whose sampling period degrades proportionally to the concurrent compute load; the dual-core and hardware-interrupt mechanisms each decouple sampling from that load in their own way. The dual-core mechanism proved the most precise in terms of effectively achieved sampling period; the interrupt mechanism offers the most exact timing reference at the trigger level, though its processing can lag under very long compute loads. Repeating the captures with more samples per run, and extending the comparison to an Arduino Uno implementation of the same mechanisms, remain as direct follow-up work for this report.
